\documentclass[11pt]{article}
\usepackage[margin=1in]{geometry}
\usepackage{amsmath,amssymb,booktabs,hyperref}

\title{Calibration of the Three-Pair Maximum-Entropy Predictor}
\author{Working explanation}
\date{5 August 2026}

\begin{document}
\maketitle

\section{What calibration computes}

Let $Y$ be the token to predict and let $A,B$ be two context tokens.
At one checkpoint the layered estimator supplies three pairwise
distributions
\[
 \widehat p_{YA}(y,a),\qquad
 \widehat p_{YB}(y,b),\qquad
 \widehat p_{AB}(a,b).
\]
The purpose of calibration is to find the maximum-entropy joint
distribution whose three pair margins equal these estimates.  Fixing the
context-pair margin allows us to write
\begin{equation}
 q_\theta(y,a,b)
 =\widehat p_{AB}(a,b)q_\theta(y\mid a,b),
 \label{eq:joint}
\end{equation}
with conditional distribution
\begin{equation}
 q_\theta(y\mid a,b)
 =\frac{p_0(y)\exp\{c_1(y,a)+c_2(y,b)\}}{Z_{ab}},
 \qquad
 Z_{ab}=\sum_zp_0(z)\exp\{c_1(z,a)+c_2(z,b)\}.
 \label{eq:conditional}
\end{equation}
Here $p_0$ is a normalized target baseline, while $c_1$ and $c_2$ are
the unknown log correction factors.  They must satisfy
\begin{align}
 \sum_b\widehat p_{AB}(a,b)q_\theta(y\mid a,b)
 &=\widehat p_{YA}(y,a),
 \label{eq:ya-margin}\\
 \sum_a\widehat p_{AB}(a,b)q_\theta(y\mid a,b)
 &=\widehat p_{YB}(y,b).
 \label{eq:yb-margin}
\end{align}
These equations explain why the third pair matters.  It supplies the
actual weights with which the two contexts co-occur when the fitted
conditional is averaged back to its target--context margins.

\section{Why the factors require a numerical solution}

Every correction changes the normalizers of many context pairs.  Therefore
the factors cannot be obtained by independently inserting
$\widehat p(y\mid a)$ and $\widehat p(y\mid b)$ into a closed formula.
They are the minimizer of the convex dual
\begin{align}
 F(\theta)
 &=\sum_{a,b}\widehat p_{AB}(a,b)\log Z_{ab}
   -\sum_y\widehat p_Y(y)\ell_y \notag\\
 &\quad-\sum_{y,a}\widehat p_{YA}(y,a)c_1(y,a)
       -\sum_{y,b}\widehat p_{YB}(y,b)c_2(y,b),
 \label{eq:dual}
\end{align}
where $p_0$ is the normalized exponential of the baseline parameters
$\ell_y$.  The gradient is the discrepancy between the margins produced
by the current factors and the desired margins:
\[
 \nabla F(\theta)=
 \begin{pmatrix}
  q_Y-\widehat p_Y\\
  q_{YA}-\widehat p_{YA}\\
  q_{YB}-\widehat p_{YB}
 \end{pmatrix}.
\]
Thus optimization of the dual and solution of the margin equations are
the same problem.

At the final $V=4096$ checkpoint the current sparse problem has roughly
$3.7$ million natural parameters: $4096$ baseline values,
$1{,}399{,}764$ first-lag corrections, and $2{,}288{,}630$ second-lag
corrections.  Constructing a Hessian is out of the question, but even one
objective-and-gradient evaluation is expensive because all observed
context edges must be normalized and accumulated into three margins.

\section{Sparse intersection evaluation}

A literal evaluation of $Z_{ab}$ sums over all $V$ targets for every
observed context pair.  To avoid this, define the deviations of the
multiplicative corrections from one:
\begin{equation}
 r_1(y,a)=\exp\{c_1(y,a)\}-1,
 \qquad
 r_2(y,b)=\exp\{c_2(y,b)\}-1.
 \label{eq:r}
\end{equation}
Define also
\begin{equation}
 S_{1a}=\sum_y p_0(y)r_1(y,a),
 \qquad
 S_{2b}=\sum_y p_0(y)r_2(y,b).
 \label{eq:S}
\end{equation}
Since $p_0$ is normalized, expanding the product in
\eqref{eq:conditional} gives
\begin{align}
 Z_{ab}
 &=\sum_y p_0(y)(1+r_1(y,a))(1+r_2(y,b))\notag\\
 &=1+S_{1a}+S_{2b}
   +\sum_y p_0(y)r_1(y,a)r_2(y,b).
 \label{eq:intersection}
\end{align}
The final sum is nonzero only where the two sparse correction rows share
a target.  We therefore precompute an intersection plan containing
\[
 (\text{context-edge index},y,c_1\text{ index},c_2\text{ index})
\]
for every nonempty intersection.  This plan produces all normalizers and
the fitted $Y,YA,YB$ margins without a dense pair or triple table.

The identity is exact mathematically but can suffer floating-point
cancellation because $r_i$ may be negative and approach $-1$.  Extreme
trial steps may therefore be unstable even though the original positive
sum in \eqref{eq:conditional} is well defined.  The current solver uses a
trust region and a direct sparse-union evaluator for small problems, but a
fully safeguarded large-problem evaluation remains necessary.

\section{Initialization and warm starts}

The first fit is initialized from the simpler pairwise conditionals
$\widehat p(y\mid a)$ and $\widehat p(y\mid b)$.  In natural-parameter
coordinates this gives a pair-product starting point, which is normally
closer than a unigram initialization.

Subsequent checkpoints are divided into two interleaved chains.  With eight
checkpoints these are
\[
 0\longrightarrow2\longrightarrow4\longrightarrow6,
 \qquad
 1\longrightarrow3\longrightarrow5\longrightarrow7.
\]
For the next problem, factors whose feature keys already existed are copied
from the preceding member of the same chain.  Newly observed features begin
at zero.  Only compact feature keys and fitted factors need to survive; the
complete preceding checkpoint can be released.

Interleaving is a compromise.  It exposes two independent fits to the CPU
while preserving warm starts.  Creating many more chains would raise
nominal parallelism but also make each chain colder, increasing the number
of iterations and potentially losing the wall-clock benefit.

\section{Current hybrid solver}

\subsection{L-BFGS approach phase}

The first phase uses gauge-reduced L-BFGS on the dual
\eqref{eq:dual}.  Global constants and redundant constants in saturated
rows are anchored.  A finite displacement region around the warm start
prevents a line-search proposal from changing log factors by hundreds.
Within that region the line search forms the adaptive mixture
\[
 \theta_{\rm trial}
 =(1-\alpha)\theta_{\rm old}+\alpha\theta_{\rm proposal},
 \qquad 0<\alpha\leq1.
\]
After an accepted iterate, the solver evaluates the actual grouped margin
certificate
\[
 C(\theta)=\max\left\{
 \lVert q_Y-\widehat p_Y\rVert_1,
 \lVert q_{YA}-\widehat p_{YA}\rVert_1,
 \lVert q_{YB}-\widehat p_{YB}\rVert_1
 \right\}.
\]
It stops once $C(\theta)$ reaches the requested tolerance.  This is more
appropriate than relying on the optimizer's largest-gradient-component
test.

\subsection{IPF polishing phase}

If L-BFGS does not reach the desired certificate, iterative proportional
fitting alternately updates the three factor families.  Schematically,
\[
 \ell_y\leftarrow\ell_y+log\frac{\widehat p_Y(y)}{q_Y(y)},
\]
\[
 c_1(y,a)\leftarrow c_1(y,a)
   +\log\frac{\widehat p_{YA}(y,a)}{q_{YA}(y,a)},
\]
and analogously for $c_2$.  Grouped inactive cells require an additional
row shift so that their aggregate mass is also correct.

One IPF sweep currently performs several complete margin evaluations: one
before the baseline update, one before each correction-family update, and
one for the final certificate.  This makes a stable but long IPF tail very
expensive.

Neither the chosen maximum of three $L^1$ errors nor every intermediate
numerical path is monotone.  The implementation therefore retains the
iterate with the best measured certificate and hands that state between
phases.

\section{Where the time goes}

The provisional full-text8 $V=4096$ run used eight checkpoints.  Its
construction took $811$ seconds, while calibration took $4801$ seconds out
of $5621$ seconds total.  Most large checkpoints required about $1100$
complete margin evaluations.

Each evaluation streams a large intersection plan and performs several
segmented reductions.  Arithmetic per entry is small, but memory traffic is
large.  At the largest checkpoints the workers compete for memory bandwidth
and allocate large partial arrays.  This explains the observed transition
from broad high-CPU construction plateaus to a long calibration phase with
only roughly three or four effective Python cores and substantial kernel
activity.  Merely adding threads will not remove this bottleneck.

\section{Present parallelism and its limitations}

Calibration is already parallelized at two levels.  First, the checkpoints
are split into independent interleaved warm-start chains.  In the provisional
$V=4096$ run, two chains were fitted concurrently.  Checkpoints within one
chain cannot be fitted concurrently in the present scheme, because checkpoint
$j+2$ is initialized with the fitted factors from checkpoint $j$.

Second, one evaluation of the margins can split the sparse intersection plan
among worker threads.  In the $V=4096$ run each fit requested five margin
workers.  The workers independently form intersection contributions and
parts of the $Y$, $YA$, and $YB$ reductions.  Thus the run could have two
simultaneous fits, each using up to five workers in its parallel portions.
The large NumPy kernels execute outside Python's global interpreter lock, so
these threads can perform useful work concurrently.

This does not mean that the complete calibration uses ten cores continuously.
Several important parts remain serial:
\begin{itemize}
\item L-BFGS chooses each trial point from the result of the preceding
      objective and gradient evaluation.  Its line search is likewise an
      adaptive sequence, so evaluations within one solve are not independent.
\item The $Y$, $YA$, and $YB$ IPF projections must be applied successively;
      a later projection uses the margins produced after the earlier update.
\item Some baseline terms, normalizers, assembly of worker results,
      diagnostics, and factor updates are still computed by the coordinating
      thread.
\item In the factorized evaluator every worker presently returns large dense
      partial reduction arrays.  Their addition is serial and their creation
      multiplies memory traffic.
\item Each warm-start chain eventually waits for its slowest large checkpoint,
      and unequal iteration counts cause load imbalance between chains.
\end{itemize}

There are two distinct obstacles to further scaling.  The first is
\emph{algorithmic dependency}: neither consecutive optimizer iterations nor
consecutive IPF projections can simply be evaluated on separate cores and
then averaged without changing the method.  Parallel speculative trial
steps are possible, but only one accepted step advances the solver, so their
benefit is uncertain.  The second is \emph{memory bandwidth}: a margin
evaluation performs relatively little arithmetic per byte read.  Once several
workers saturate the memory subsystem, more cores mainly compete for the same
data and produce more temporary arrays rather than reducing elapsed time.

The most credible additional parallelism is therefore structural rather than
merely increasing the thread count.  Workers should own disjoint output
ranges and write compact partial results, avoiding full-sized replicated
reductions and a serial merge.  Independent chains, experiments, or
speculative solver starts can occupy otherwise idle cores.  More importantly,
incremental IPF or a better-preconditioned second-order method would reduce
the number of complete margin scans; this attacks the dominant cost even when
memory bandwidth prevents a large speedup from additional cores.

\section{Most promising acceleration routes}

The largest opportunity is to reduce the number of global passes.

\begin{enumerate}
\item \textbf{Closer continuation.}  Thirty-two checkpoints make adjacent
      problems more similar.  Better warm starts may substantially reduce
      iterations, so 32 checkpoints need not cost four times as much as
      eight.  Preserving quasi-Newton history or a compact preconditioner
      across checkpoints should also be investigated.

\item \textbf{Progressive vocabulary.}  Fit successively at
      $V=256,512,1024,4096,$ and the full vocabulary, transferring factors
      at each expansion.  This avoids fitting a huge alphabet from a tiny
      prefix while allowing calibrated prediction to begin early.

\item \textbf{Incremental IPF.}  A first-lag update for context $a$ affects
      only $AB$ edges incident to $a$; a second-lag update for $b$ affects
      only edges incident to $b$.  Maintaining affected normalizers and
      partial margins could replace several full scans by local graph
      updates.

\item \textbf{Stronger second-order methods.}  The Hessian is a covariance
      operator for the sufficient statistics.  It should not be stored, but
      Hessian-vector products, Newton--CG, block-diagonal preconditioners, or
      the two-pair tree solution as a preconditioner may reduce hundreds of
      passes to tens.

\item \textbf{Output-sharded reductions.}  Each worker should own disjoint
      YA/YB output ranges instead of returning complete partial arrays.
      This reduces memory traffic and is important for the 64-core,
      240GB server.

\item \textbf{Blockwise intersection plans.}  Streaming or memory-mapping
      plan blocks can reduce peak memory.  If necessary, one complete
      checkpoint can be fitted at a time while compact warm starts for both
      logical chains are retained.

\item \textbf{Precomputation.}  With a fixed corpus, alphabet, and 32
      checkpoint boundaries, cache count increments, layered profiles and
      their evaluations, and optionally prepared sparse checkpoint problems.
      This does not eliminate calibration, but removes repeated construction
      and makes solver experiments and restarts inexpensive.
\end{enumerate}

\section{Additional acceleration ideas retained for investigation}

The following proposals are sufficiently well founded to keep on record.
They are possible changes to the numerical route, not changes to the
maximum-entropy model or to the final margin certificate.

\subsection{Stochastic approach phase}

The dual gradient is an expectation over observed context pairs
$(a,b)\sim\widehat p_{AB}$.  Sampling edges with this probability therefore
gives an unbiased stochastic gradient (or, with a different sampling law, an
importance-weighted unbiased gradient).  Far from the solution, Adam, SVRG,
or another stochastic method could use one to five percent of the edges per
step.  Occasional exact evaluations would monitor progress, and an exact
IPF, L-BFGS, or Newton phase would still produce the final certificate.

This is promising because it attacks the roughly $1100$ complete passes,
but it must be tested carefully.  Rare $YA$ and $YB$ features may otherwise
be visited too infrequently, so stratification by context or feature may be
needed.  A subsample cannot certify convergence, and normalization over the
target support is still required for every sampled edge.

The first implementation uses contiguous probability-weighted edge blocks.
A block is chosen with probability equal to its total $\widehat p_{AB}$ mass
and is evaluated with the conditional edge distribution inside the block.
Its gradient is therefore unbiased.  The block is a view into the permanent
intersection plan rather than a copied random subplan, and independent blocks
can be evaluated concurrently.  An SVRG control variate uses
\[
 \widehat g_b(\theta)-\widehat g_b(\widetilde\theta)
 +\nabla F(\widetilde\theta),
\]
where $\widetilde\theta$ is an occasional exactly evaluated snapshot.  This
remains unbiased and greatly reduces block-to-block variation near a warm
start.

In a provisional final-checkpoint $V=4096$ probe, one full gradient took
$16.14$ seconds on one CPU worker, while a typical zero-copy block gradient
took about $0.17$ seconds.  Twenty-five four-replica SVRG updates reduced the
exact gradient $\ell_2$ norm from $0.003472$ to $0.002037$.  Counting both the
current and snapshot block evaluations, they touched about $0.90$ million
edge instances, versus $1.40$ million edges in one full pass.  These are
engineering diagnostics, not final solver timings: exact snapshot passes,
plan construction, and the subsequent certified handoff must all be included
in an end-to-end comparison.

A subsequent controlled $V=1024$ last-checkpoint experiment reached the
actual grouped certificate $0.00912<0.01$ using SVRG plus occasional exact
certificate evaluations and no old L-BFGS/IPF finishing phase.  The
stochastic phases took about $279$ seconds in several restartable runs,
compared with $674$ seconds for the original exact fit of the same checkpoint,
a provisional $2.4\times$ elapsed-time improvement.  Splitting the run caused
the fixed plan to be rebuilt more than once, so an integrated run should be
somewhat faster.  In contrast, handing an insufficiently converged stochastic
state to the old solver saved only about one percent and is not the proposed
route.

The statistical batch size and the execution worker count must be separate
parameters.  An initial comparison inadvertently increased both together and
therefore changed the optimizer.  In the corrected scaling test, every update
used the same twelve sampled blocks and produced bit-identical factors and
certificate traces.  Four workers took $42.92$ seconds, eight took $37.91$
seconds, and twelve took $35.70$ seconds.  Twelve workers are therefore best
for this fixed computation on the 15-core laptop, although the modest
$1.20\times$ gain over four workers shows substantial memory-bandwidth
contention.  Batch size remains a separate optimization question.  Both
counts are machine-dependent and must be re-benchmarked on the server;
elapsed time to the identical certificate, not utilization or nominal
parallelism, is the selection criterion.

\subsection{Context-owned IPF and two incidence orders}

Within a first-factor projection, the factor rows indexed by distinct
contexts $a$ can be owned by distinct workers; the analogous statement holds
for the $b$ rows in the second-factor projection.  Store the static incidence
plan in two CSR-like orders, one grouped by $a$ and one by $b$.  A worker then
writes only its owned factor and margin ranges.  This can remove replicated
full-sized partial arrays and their serial merge.  Shared target-margin
reductions and the succession of the three projection families remain.

Fusing the normalizer and margin calculations, or maintaining affected
normalizers incrementally, may also reduce the number of complete evaluations
per IPF sweep.  Any such implementation must be checked against the present
Gauss--Seidel IPF map; reusing stale margins silently changes the iteration.

\subsection{Fixed-point and second-order acceleration}

Guarded overrelaxation is cheap to test.  Anderson acceleration is already
implemented in the code base and should be reassessed on the present large
sparse problems rather than assumed to provide a particular gain.  Late in
the solve, Greenkhorn-style updates of only the worst context rows could turn
global sweeps into local graph updates.  A matrix-free Newton--CG endgame is
also plausible: Hessian--vector products use the same covariance structure as
the gradient, while a two-pair tree solution or rowwise Fisher blocks could
serve as a preconditioner.

Newton--CG is deferred, not discarded.  A transparent matrix-free
Hessian--vector product has been derived by differentiating the sparse
factorized margin algebra and checked against centered finite differences;
it also has the required positive-semidefinite and gauge-null properties on
the tests.  This gives a precise starting point if a second-order solver is
revisited.  The current development priority remains the parallel stochastic
solver with a standard learning-rate scheduler.

\subsection{Continuation across checkpoints}

A single chain uses the immediately preceding checkpoint and requires only
one cold start.  Once feature supports have been aligned and a common gauge
has been fixed, a guarded secant extrapolation from the previous two fits may
be closer still.  Starting later checkpoints from unfinished partial iterates
is recorded as a more speculative option: it increases concurrency, but also
duplicates large states, consumes bandwidth, and may discard work when the
predecessor subsequently moves.

\subsection{Compact and mixed-precision evaluation}

Use 32-bit plan indices wherever their range is sufficient, make one index
implicit through CSR offsets, and consider delta encoding within sorted runs.
Plan values and factors can be tested in \texttt{float32}, while sensitive
normalizers, reductions, objectives, and the final certificate remain in
\texttt{float64}.  Agreement with the all-double reference, especially near
the stopping tolerance and on extreme trial steps, is required.

The fast factorized expansion should be used only where it is numerically
safe.  Shards with corrections near the cancellation boundary can fall back
to the positive sparse-union log-sum-exp evaluation.  Accelerated hardware
makes the safe path cheaper but does not remove this numerical requirement.

\subsection{Sparse kernels and the local Apple silicon}

The intersection calculation can be viewed as a sampled weighted product,
and the margin calculations as sampled products followed by sparse
reductions.  This vocabulary may help in selecting optimized kernels, but it
does not by itself replace the current intersection algorithm: the two factor
matrices are sparse, whereas many standard sampled dense--dense matrix
multiplication interfaces assume dense operands.

The present laptop reports an Apple M5 Pro with 15 CPU cores, a 16-core GPU,
and 24 GB of unified memory.  Apple's specification gives this configuration
307 GB/s unified-memory bandwidth and Neural Accelerators in the GPU cores.
Metal Performance Shaders supplies optimized matrix multiplication, and
Apple's newer Metal Performance Primitives expose tensor operations intended
to use the M5 GPU's neural accelerators.  These facilities make a local GPU
kernel experiment worthwhile.  They do not guarantee acceleration of our
irregular sparse gather and segmented-reduction workload, nor is the separate
Neural Engine a general-purpose target for arbitrary solver code.  The first
step is therefore a small benchmark of one fixed saved-checkpoint margin
evaluation, including conversion, precision, memory use, and exact agreement,
before any solver port.

\subsection{Measured stochastic continuation and the remaining memory peak}

The first full-file test of this strategy used $V=1024$ and 32 geometric
checkpoints.  Checkpoint problems were constructed once and persisted; a
single calibration chain then used the pair-product start at the first
checkpoint and the preceding certified solution thereafter.  Each fit used
12 parallel block-SVRG samples per Adam update, an exact certificate every
50 updates, and a standard reduce-on-plateau learning-rate schedule.  An
exact factorized L-BFGS fallback was available at every checkpoint, starting
from the best stochastic candidate.

All 32 checkpoints reached grouped margin certificate $0.01$ without exact
fallback or nonfinite rejection.  Calibration took $142.5$ seconds in total:
$96.4$ seconds in sampled gradients, $15.7$ seconds in reference-block
caches, $9.5$ seconds in optimizer updates, and $6.0$ seconds in exact
certificate evaluations.  This validates the stochastic solver as the
production first attempt, while retaining exact fitting as a correctness
backstop.

The $10.25$ GB peak resident set is not an accumulation of checkpoint
states.  The final saved problem occupies only $29$ MB, and only the current
and immediately preceding problems coexist during warm-start transfer.  The
dominant object is instead the final checkpoint's $83{,}702{,}774$-entry
intersection plan: its four 32-bit index arrays occupy $1.34$ GB, while the
temporary SciPy sparse products used to construct it raise resident memory
to $7.78$ GB before stochastic caches and concurrent outputs are counted.
The next implementation step is therefore a direct or streamed construction
of those compact arrays, avoiding the sparse-matrix temporaries.  This change
should increase the alphabet attainable on a 24 GB laptop without changing
the estimator or the solver trajectory.

That direct construction is now implemented in the native extension.  It
forms context-row indices for both factors, merge-intersects their sorted
target lists for each observed edge, counts the matches in one pass, and
allocates and fills only the four final 32-bit arrays in a second pass.  On an
intermediate saved $V=1024$ checkpoint with $4{,}332{,}210$ intersections,
the sparse-product reference required $0.462$ seconds and peaked at $1.13$
GB, whereas the direct builder required $0.185$ seconds and peaked at $160$
MB; the final plan itself was $69.3$ MB.  Randomized tests require exact
array equality with the retained SciPy reference.

The SVRG reference cache has also been made context-sparse.  A block retains
only the YA and YB correction positions whose contexts occur on its edges,
rather than two full correction vectors, and uses 32-bit positions.  Current
RSS and the exact byte sizes of the plan and cache are now recorded at every
checkpoint.  Finally, a plateau at the minimum learning rate terminates the
stochastic phase and transfers its best candidate to the exact fallback,
rather than exhausting an arbitrary maximum update count.

A complete rerun of the same 32-checkpoint chain validates these memory
changes without a statistical or numerical confound.  All persisted fitted
arrays, update counts, and certificates agree exactly with the earlier run;
there are again no exact fallbacks.  Elapsed calibration time is $151.3$
seconds rather than $142.5$ seconds, while measured peak RSS falls from
$10.25$ GB to $5.12$ GB.  At the final checkpoint the direct intersection
plan occupies $1.339$ GB and the context-sparse reference cache $0.776$ GB;
parent RSS after releasing the active fit is $2.72$ GB.  The remaining gap
between live-array sizes and RSS is allocator-retained memory, not retained
checkpoint states.

Predictive scoring retains the positive sparse-union log-sum-exp calculation.
A faster factorized-normalizer experiment was rejected because cancellation
changed the aggregate result slightly.  Instead the 31 independent
checkpoint intervals are distributed across processes that share a
memory-mapped reduced token stream.  Four workers reproduce every serial
interval and every honest-code accounting total exactly, reducing elapsed
scoring time from $256.8$ to $87.7$ seconds ($2.93$-fold).

The subsequent full-file $V=4096$, 32-checkpoint calibration also completed
without fallback or nonfinite rejection.  Its 20,600 stochastic updates
took $1330.6$ seconds and peaked at $13.89$ GB including workers.  The final
plan and cache occupy $4.48$ and $3.19$ GB respectively.  Exact scoring
revealed a load-balancing issue because intervals grow over time; scheduling
the longest intervals first and restoring causal order before summation
reduced four-worker time from $647.7$ to $553.6$ seconds.  Increasing to 8
and 12 workers reduced this further to $321.6$ and $270.8$ seconds.  All
rows and totals are exactly identical, so the additional utilization is
useful reference computation rather than speculative duplicated fitting.

\section{Conclusion}

The dominant task is a large convex moment-matching problem: find the
natural parameters of a conditional exponential family so that averaging
against the observed context-pair distribution reproduces two target--lag
margins.  The sparse intersection identity makes each global pass possible,
and warm starts plus hybrid L-BFGS/IPF make the solve reliable.  The current
limitation is the product of two quantities: too many global passes and too
much memory traffic per pass.  Reducing the pass count, exploiting local
updates, and partitioning reductions by owned outputs are therefore the
main routes to a full-vocabulary, 32-checkpoint calculation.

\end{document}
